% ===========================================================================
\chapter{Asking first}\label{ch:permission}

Our agent can now read, search, edit and run anything in its workspace.
Who decides whether it should? So far, \code{Allow_Everything}. That is
fine for a scratch copy you are happy to throw away, and that is what
\code{--yes} will mean in Chapter~\ref{ch:program}. For anything else, the
agent should ask you first.

The line is drawn by the tools themselves: each says whether running it
\code{changes_things}. Reading and searching don't, and are always allowed;
asking you every time the agent wants to look at a file would make it
useless, and looking can't do harm in a workspace that holds nothing
secret. Editing, writing and running commands do change things, and need
a yes.

Every permission is only asked about calls whose input has been checked
(that is \code{Permission}'s precondition, from Chapter~\ref{ch:loop}), so
a policy that shows you the call can rely on its arguments being there.

\section{Read-only}

The simplest policy that says no allows only tools that change nothing.
With it, the agent can still explore a project and explain it, or find a
bug and tell you where it is, without being able to touch anything:

\begin{append}{permissions.nex}
-- Only tools that change nothing: `--read-only`. The agent can still
-- explore the project and explain it.
class Allow_Reading inherit Permission
create
  make() do
  end

feature
  allows(tool: Tool, input: Json_Value): Boolean do
    result := not tool.changes_things()
  end
end
\end{append}

\section{Asking you}

\code{Ask_First} asks at the keyboard, about the call the agent has just
printed (\code{run_command nex checks.nex}, say). For an edit it also shows
the lines going out and coming in, the way a diff does, so you know what
you are allowing.
Besides yes and no, you can answer ``always'', which allows that tool for
the rest of the session: after you have seen the agent make a few good
edits, approving each one becomes a chore.

\begin{append}{permissions.nex}
-- Ask the person at the keyboard before anything that changes things.
-- "a" (always) allows that tool for the rest of the session.
class Ask_First inherit Permission
create
  make() do
  end

feature
  always: Array[String]    -- tools the user has allowed for good

  allows(tool: Tool, input: Json_Value): Boolean do
    if not tool.changes_things() or always.contains(tool.name()) then
      result := true
    else
      result := ask(tool, input)
    end
  end

  private feature
    -- Ask about the call the agent has just printed.
    ask(tool: Tool, input: Json_Value): Boolean do
      if tool.name() = "edit_file" then
        console.print(indented(edit_preview(input), "      "))
      end
      let answer: ?String := console.read_line("    allow? [y]es, [n]o, [a]lways for "
        + tool.name() + "> ")
      if ?answer as text then
        let a: String := text.trim().to_lower()
        if a = "a" or a = "always" then
          always.add(tool.name())
          result := true
        else
          result := a = "y" or a = "yes"
        end
      else
        result := false     -- no answer (end of input) is a no
      end
    end

    -- An edit as lines taken out and lines put in.
    edit_preview(input: Json_Value): String do
      result := indented(input.get("old_text").as_string, "- ")
        + indented(input.get("new_text").as_string, "+ ")
    end
end
\end{append}

\code{console.read_line} gives a \code{?String}: \code{nil} at the end of
the input (Ctrl-D, or a program whose input has run out). Anything but a
clear yes is a no.

\section{Answers given in advance}

Checks can't type at a keyboard, so there is one more policy, for them: a
list of answers, given in advance, used in order.

\begin{append}{permissions.nex}
-- Answers given in advance, in order, for checks: true allows the next
-- tool that asks, false refuses it.
class Answers_Given inherit Permission
create
  make(answers: Array[Boolean]) do
    this.answers := answers
  end

feature
  answers: Array[Boolean]
  asked: Integer

  allows(tool: Tool, input: Json_Value): Boolean do
    if not tool.changes_things() then
      result := true
    else
      result := asked < answers.length and answers.get(asked)
      asked := asked + 1
    end
  end
end
\end{append}

Here is a script that tries the fix twice. The first time it is refused;
the second time it is allowed. Watch what the model is told:

\begin{tryit}{try_permission.nex}
intern Agent

let ws: Workspace := fresh_workspace(user_path("sample_project"), user_path("scratch"))
let fix: String := '{"path": "loans.nex", "old_text": "    if day >= due_on() then\n      result := day - due_on() + 1", "new_text": "    if day > due_on() then\n      result := day - due_on()"}'
let model: Scripted_Model := create Scripted_Model.make([
  use_tool("A book is late only after its due day.", "edit_file", fix),
  use_tool("May I? It is a two-line change to days_late.", "edit_file", fix),
  say("Done.")
])
let permission: Answers_Given := create Answers_Given.make([false, true])
let agent: Agent := create Agent.make(model, standard_tools(ws), permission,
  create Reporter.make(false, false))
agent.run("Fix the late fines.")
console.print_line("asked " + permission.asked.to_string + " times")
\end{tryit}
\tryoutput{try_permission}

The refusal reached the model as an error, with the request not to look
for another way round it. You will see \code{Ask_First} at work in
Chapter~\ref{ch:program}, once there is a program to type at.
